\section*{Bab \ref{ch:g1:shapes} --- Bangun Datar dan Letak}

\begin{solution}{exo:g1:shapes:1}
$3$ titik sudut: segitiga. Bundar, tanpa titik sudut: lingkaran. $4$
sisi sama panjang: persegi.
\end{solution}

\begin{solution}{exo:g1:shapes:2}
Jawaban beragam: pintu dan buku berbentuk persegi panjang; piring
atau jam dinding berbentuk lingkaran; rambu peringatan jalan
berbentuk segitiga; ubin bisa berbentuk persegi.
\end{solution}

\begin{solution}{exo:g1:shapes:3}
Salah: persegi yang diputar sampai berdiri di atas titik sudutnya
tetap punya empat sisi sama panjang dan empat titik sudut ---
memutar tidak mengubah bangunnya.
\end{solution}

\begin{solution}{exo:g1:shapes:4}
Gambar di kertas berpetak: persegi berukuran $3 \times 3$ petak,
persegi panjang $5 \times 2$, dan segitiga memakai garis petak yang
miring sebagai salah satu sisinya.
\end{solution}

\begin{solution}{exo:g1:shapes:5}
Jawaban beragam. Contoh: bola ada di sebelah kiri, mobil ada di
antara bola dan boneka, boneka ada di sebelah kanan.
\end{solution}

\begin{solution}{exo:g1:shapes:6}
Tanda silang digambar di dalam lingkaran, bintang di luar, dan titik
tepat pada garis lingkarannya.
\end{solution}

\begin{solution}{exo:g1:shapes:7}
Mengikuti $\to\ \uparrow\ \to\ \uparrow\ \to$ berarti bergerak $3$
petak ke kanan dan $2$ petak ke atas. Jalan pulangnya adalah
$\downarrow\ \downarrow$ lalu $\leftarrow\ \leftarrow\ \leftarrow$
(atau jalur mana pun yang kembali ke petak awal).
\end{solution}

\begin{solution}{exo:g1:shapes:8}
Dua segitiga dan satu persegi: $3 + 3 + 4 = 10$ titik sudut. Satu
lingkaran dan dua persegi panjang: $0 + 4 + 4 = 8$ titik sudut.
\end{solution}

\begin{solution}{exo:g1:shapes:9}
$12$ batang korek api membentuk persegi dengan $3$ batang di setiap
sisi ($3 + 3 + 3 + 3 = 12$). Dengan $10$ batang tidak bisa: $10$
tidak terbagi menjadi $4$ bagian utuh yang sama (setiap sisi akan
memerlukan $2$ setengah batang).

\end{solution}
